%!TEX root = rapport.tex

\section{Preliminaries: Basic Lambda Calculus}

The untyped lambda calculus was introduced in 1936 by Alonzo Church
[reference] as part of the search for a solid foundation for mathematics. It is
heavily intertwined with the notion of a ``function'', in the sense that logic
and mathematics is represented with only functions and untyped variables that
can be either free or bound.


% alt er funktioner
% eksempler

As a very informal example, here is a function that doubles an integer:
\begin{equation*}
	(\lam{x}{x+x})2=4
\end{equation*}
This example shows a bound variable, $x$, and an implicit \emph{reduction} step. The variable $x$
is said to be bound because the lambda function uses it as an argument. Contrast this to
\begin{equation*}
	(\lam{x}{x+y})2=2+y.
\end{equation*}
Here, $y$ is a free variable that hasn't got any values assigned to it yet. It can
take any value, hence it's called ``free''. In this second example there was
another reduction step. A reduction step can be regarded as a rewriting action
on the string representing the lambda expression. Yet more informally, one can 
think of it as ``substituting in'' the values. In the examples above this amounts to
replacing $x$ with the number 2. 

% definition af syntax

The lambda calculus consists of \emph{lambda expressions} and variables with a
recursively defined syntax. The set of all lambda expressions is denoted $\Lambda$:
\begin{align*}
	x \text{ is a variable} &=> x\in\Lambda \\
	x \text{ is a variable and } M\in\Lambda &=> (\lam{x}{M})\in\Lambda \\
	M, N\in\Lambda &=> (M N)\in\Lambda
\end{align*}
For brevity, it's convention to leave out the outermost parentheses.
Expressions are evaulated left-to-right, so $M N P = (M N) P$, and function
bodies extend as far as possible; $\lam{x}{M N}=\lam{x}{(M
N)}\neq(\lam{x}{M})N$.

% frie variable
As already indicated, the variabled in a lambda expression can be split on two groups;
the free and the bound variables. The function $FV$ takes returns for any
lambda expression the set of free variables:
\begin{align*}
	FV(x)&={x} \\
	FV(\lam{x}{M})&= FV(M) \backslash {x} \\
	FV((M N)) &= FV(M) \cup FV(N)
\end{align*}

% definition af alfakonvertering
% eksempel
% definition af betareduktion
% eksempel!
% definition af reduktionsgraf
% eksempel(?)


% Syntax

% Alpha-conversion

% Beta-reduction